\section{Introduction}
\label{sec:intro}

Micropollutants, including pharmaceuticals, pesticides, endocrine
disruptors, and personal-care products, are detected at ng/L to \textmu g/L
levels in surface water, groundwater, and treated effluent worldwide
\citep{TODO-micropollutant-occurrence}. Conventional treatment trains
remove many of these compounds only partially, and their chronic ecological
and human-health effects have placed micropollutant control on regulatory
agendas in the European Union, the United States, and Asia
\citep{TODO-regulation}. Adsorption is among the most widely deployed
polishing steps because it requires no reagent input, produces no
transformation by-products, and operates across wide concentration ranges
\citep{TODO-adsorption-review}.

Layered double hydroxides (LDHs) are a class of anionic clays with the
general formula
$[\mathrm{M}^{2+}_{1-x}\mathrm{M}^{3+}_{x}(\mathrm{OH})_2]^{x+}
[\mathrm{A}^{n-}]_{x/n}\cdot m\mathrm{H_2O}$, whose metal composition,
interlayer anion, and calcination state can be tuned to target specific
pollutants \citep{TODO-ldh-review}. Hundreds of LDH variants have been
synthesized and tested against individual micropollutants, and each such
study reports removal under its own grid of pH, temperature, contact time,
adsorbent dose, and initial concentration. The resulting literature is
large but fragmented: no single laboratory can cover the combinatorial
space of metal pairs, synthesis routes, and pollutants experimentally.

Isotherm and kinetic models describe each of these systems in isolation
and offer no mechanism for pooling evidence across studies. Machine
learning (ML) offers a natural response to this fragmentation.
By compiling published experiments into a single dataset, ML models have
been trained to predict adsorption capacity or removal ratio from
material, pollutant, and operating descriptors, with reported test
$R^2$ values frequently above 0.9
\citep{TODO-ml-adsorption-1, TODO-ml-adsorption-2, TODO-biochar-cubist}.
Such models promise two distinct services: interpolating the behavior of
an already-studied adsorbent--pollutant system at untested conditions, and
screening new systems before synthesis. The second service is the one
most frequently invoked to motivate these studies.

However, we argue that the evaluation protocols behind these reports do
not support the screening claim. Firstly, literature-compiled datasets
are strongly clustered: in our compilation, 1,342 experimental records
originate from only 59 source studies, and within one study the material
and pollutant descriptors are constant while only operating conditions
vary. A random record-level train--test split therefore places
near-duplicate records on both sides, and the test score measures
within-study interpolation while being reported as generalization.
Secondly, preprocessing steps such as imputation, scaling, and feature
selection are commonly fitted on the full dataset before splitting, which
leaks test-set information into training \citep{kapoor2023leakage}.
Thirdly, point predictions are reported without uncertainty estimates or
applicability bounds, leaving practitioners no way to know when a
prediction should be trusted.

To address the above challenges, we evaluate ML models for micropollutant
adsorption on LDHs under a protocol that closes each of these leakage
pathways. In essence, our idea is
to treat the source study as the unit of generalization: we evaluate every
model twice, once under a duplicate-free record-level split that answers
the interpolation question, and once under grouped, leave-one-study-out
protocols that answer the screening question, with all preprocessing
confined to training folds in both cases. We complement point predictions
with split-conformal intervals and a distance-based applicability-domain
score, and we examine whether chemistry-informed features improve
cross-study transfer.

More specifically, this study aims to (1) quantify the gap between
interpolation and cross-study generalization for seven regression models
on 1,342 LDH--micropollutant records from 59 studies; (2) equip the
predictions with distribution-free uncertainty intervals and an
applicability-domain flag, and test their reliability under both tasks;
and (3) evaluate chemistry-informed features, including the electrostatic
driving force pH$-$pH$_{\mathrm{pzc}}$, and read the resulting feature
attributions against known adsorption mechanisms. To the best of our
knowledge, this is the first study of LDH adsorption modeling that
separates these two prediction tasks explicitly and reports what
literature-trained models can, and cannot yet, deliver for practice.
